\newpage
\thispagestyle{empty}
\chapter*{\centering Project Overview}
\addcontentsline{toc}{chapter}{Project Overview}

A practical lab exam is hard to supervise. One or two invigilators watch 30 to 40 screens, attendance is marked on paper, and code is collected on USB drives. SmartExam, the system built for our registered project ``AI-Driven Lab Exam Monitoring and Management System'', replaces this with three parts: a Windows Presentation Foundation (WPF) client on each lab PC, a React web panel for staff and students, and an ASP.NET Core 9 API backed by PostgreSQL on Neon.

Teachers create exams with multiple-choice (MCQ), coding and theory questions, set eligibility from imported attendance and assign each student a workstation. A student can log in only from the workstation bound to the account by its hardware ID. After a readiness check, the client locks the PC: fullscreen window, blocked shortcuts, clipboard guard and a check for unapproved programs. It reports heartbeats, focus and process events, screenshots and, when the exam requires it, webcam checks (camera covered or frozen, no face or several faces, head turned away, phone in view); no face is matched to an identity. A rule-based classifier rates each student, and invigilators follow every seat on a live SignalR dashboard, from which they can warn a student, extend time, unlock the PC or force a submission. Each edit goes to an encrypted local backup within a second and to the server within 10 seconds. Practice exams let students rehearse without affecting their results.

After each submission, C and C++ answers run against the teacher's test cases in an isolated Docker container where available (otherwise with a local compiler that is not isolated), theory answers get a TF-IDF similarity score, and k-gram winnowing flags similar code. The teacher reviews every suggested mark. Security-relevant actions are kept in a hash-chained, tamper-evident audit log, and teachers export PDF reports, an Excel workbook and an evidence ZIP. We built the system in eight four-week Agile sprints. Automated unit, integration, API and browser tests cover all twelve modules (Chapter~7); checks that need a real lab PC, such as the webcam, USB and 40-client load checks, are still pending.
